% Generated from closed P55 evidence; no model or task execution.
\newcommand{\InventoryAllCells}{%
\begin{tabular}{llrr}\hline 課題 & 条件・成否 & 仕事全体token & 継続生成\\\hline
写真整理 & F・未達 & 145,850 & 8\\
 & D・成功 & 120,935 & 6\\
 & B・未達 & 165,015 & 11\\
電気代精算 & F・未達 & 163,918 & 9\\
 & D・未達 & 161,225 & 9\\
 & B・未達 & 154,720 & 9\\
返金 & F・成功 & 85,697 & 2\\
 & D・成功 & 117,223 & 4\\
 & B・成功 & 117,305 & 4\\
曲一覧の書出 & F・未達 & 117,100 & 3\\
 & D・未達 & 156,265$+?$ & 7\\
 & B・未達 & 126,734 & 5\\
\hline\end{tabular}}
\newcommand{\InventoryTotals}{%
\begin{tabular}{lrrr}\hline 指標 & F & D & B\\\hline
成功／評価済み & 1/4 & 2/4 & 1/4\\
仕事全体・既知token & 512,565 & 555,648 & 563,774\\
継続・既知token & 302,185 & 345,268 & 353,394\\
使用量不明生成 & 0 & 1 & 0\\
前半＋継続生成 & 54 & 58 & 61\\
継続生成 & 22 & 26 & 29\\
前半＋継続API & 224 & 213 & 173\\
継続API & 173 & 162 & 122\\
復元API & 51 & 51 & 51\\
継続時間（秒） & 1,308.5 & 1,354.0 & 1,359.1\\
\hline\end{tabular}}
\newcommand{\InventoryFirstInputs}{%
\begin{tabular}{lrrr}\hline 課題 & F & D & B\\\hline
写真整理 & 8,789 & 7,342 & 6,848\\
電気代精算 & 11,267 & 11,009 & 10,202\\
返金 & 15,746 & 14,964 & 14,964\\
曲一覧の書出 & 13,771 & 12,142 & 9,473\\
\hline\end{tabular}}
\newcommand{\InventoryResultSummary}{全12予定条件のうち12件を評価し，0件が未評価，0件が未実行だった．継続費用は既知1,000,847 token，使用量不明は1生成だった．詳細省略Dは固定の採用基準を満たさず，品質を保った総費用削減として採用しない．Dは既知分だけでFを43,083 token上回るため，未知分を除いても削減にはならない．}
\newcommand{\InventoryContrasts}{全表示Fに対する全４課題の確定費用差は算出しない．一覧も省くBに対する全４課題の確定費用差は算出しない．}
\newcommand{\InventoryCaseEvidence}{写真整理のF/Dは，それぞれ42／42件のファイル移動要求を送ったが，成否は未達／成功だった．この対では全ての非supervisorアプリの保存ファイルが一致し，異なるアプリ状態はsupervisorだけだった．完了通知の要求はFで０，Dで１件であり，公式成否の差をファイル整理そのものの改善と呼ばない．隠れた個別テストの原因同定ではない．返金はF/D/Bの全条件が成功した．継続tokenはF/Dで32,209／63,735であり，初回の詳細省略量だけから終了までの費用を予測できない．}
\newcommand{\InventoryVariation}{過去P47と今回P55の全表示条件について，両方を評価できた比較のうち1課題で成否が異なった．これは事前に反復数を割り付けた追試ではなく，変更の原因や分散の推定には用いない．}
